\section{Introduction}
\label{sec:intro}

The diagnostic content of rolling-element bearing vibration sits far above its modulation rate.
A localized defect excites structural resonances, and the fault frequency only appears as the
repetition rate of the excited transients. In the Paderborn benchmark
\citep{lessmeier2016paderborn}, the resonance that carries an electrical-discharge outer-race
defect occupies roughly 2.5--4.4~kHz, and the one that carries inner-race defects 8.4--11~kHz,
while the fault orders are 3.05 and 4.95, that is, about 76~Hz and 124~Hz at 1500~rpm
(Section~\ref{sec:data}). Sampling such signals at twice their highest frequency is costly in
energy and in data, and it is not what the information requires.

The rate that the information requires is set by the occupied bandwidth, not by the highest
frequency. This is the content of Landau's theorem \citep{landau1967}: the minimal average
sampling density of a signal is the measure of its spectral support. Bandpass sampling realizes
it for a single band with uniform sampling at suitable rates \citep{vaughan1991}, and multicoset
sampling realizes it for unions of bands at any position
\citep{venkataramani2001optimal,mishali2009xampling}. Operating below the Nyquist rate is therefore
not a departure from sampling theory; it is the regime that sampling theory prescribes for
structured signals. Envelope analysis of bandpass-sampled vibration has been shown to keep bearing
fault signatures well below the Nyquist rate \citep{lu2018undersampled}.

Blind envelope analysis, however, needs a band selector, and the standard one is spectral kurtosis
through the fast kurtogram \citep{antoni2007kurtogram}. Spectral kurtosis measures impulsiveness,
not periodicity, and it is known to be captured by impulsive interference that is not related to
the fault: electromagnetic interference from variable-frequency drives \citep{smith2016emi,
vaerenberg2025cyclicemi} and non-cyclic impulsive noise in general
\citep{hebdasobkowicz2020nongaussian,hebdasobkowicz2025robustcorr}. All of these studies assume a
conventional acquisition, with an anti-alias filter in front of the converter. In that setting the
filter confines the interference to the acquisition band; \citet{vaerenberg2025cyclicemi} even
note that the anti-alias filter acts as the carrier band of the EMI impulses.

Pure decimation below the Nyquist rate removes that filter by design, because bandpass sampling
needs the resonance band to fold into baseband. Whatever impulsive content lies outside the band
of interest folds as well, and it may land on the range that the selector searches. Whether, when
and how strongly this happens has not been studied, as far as our searches show: keyword searches
and forward-citation tracking of \citet{smith2016emi} and \citet{vaerenberg2025cyclicemi} found no
work on band selection under decimation or folding (Section~\ref{sec:related}).

This paper studies that interaction and its consequence for learned acquisition.
\todo[inline]{Contributions: confirm each against the final results before submission.}
\begin{enumerate}
  \item We show, on the public Ottawa benchmark \citep{huang2018ottawa}, that single-sample spikes
  present in the raw records, healthy ones included, capture kurtogram band selection under pure
  decimation and produce a non-monotonic loss of order-estimation accuracy with the decimation
  factor. Removing the spikes before decimation restores the accuracy, in a pre-registered test.
  Three simple steps that need no access to the full-rate record do the same at $4\times$:
  cleaning isolated outliers in the decimated stream, a three-point median, and a lower search
  ceiling. Each has a limit, which we report.
  \item We reproduce the capture causally by injecting impulsive interference into the CWRU
  benchmark \citep{smith2015cwru}, against a phase-randomized surrogate with the same spectrum.
  \item We test the same interventions on the Paderborn benchmark, whose drive is fed by a 16~kHz
  PWM inverter. There, the kurtogram is already captured at full rate, decimation produces no
  valley, spike removal changes little, and the SES infogram \citep{antoni2016infogram} helps only
  at the lowest rate. A lower search ceiling at full rate gains more than decimation, though short
  of significance in a pre-registered test. The replication shows when the mechanism's
  precondition, a selector that is right at full rate, is absent.
  \item We show with a pre-registered power gate that a widely used benchmark cannot detect any
  improvement over a training-selected uniform rate, not even from a perfect estimator, while
  another can. We argue that such a gate should precede the training of learned acquisition on
  condition-monitoring data.
  \item We test learned sampling masks under a hardware-feasible set on both benchmarks, at equal
  or lower average rate. The learned mask does not outperform uniform decimation trained on the
  same records or a random mask at any budget from 1/4 to 1/32. On Paderborn, where a gain of 5
  points was detectable, the differences stay within $-5.2$ and $+4.4$ points, and the margin over
  the classical pipeline belongs to the decoder that all masks share.
\end{enumerate}
